% =============================================================================
% TEMPLATE DISPENSE UNIVERSITARIE — VARIANTE SCHEMATICA / OUTLINE
% Stile "da appunti": definizioni con i due punti, liste annidate,
% sotto-titoli in grassetto, alberi di classificazione (forest).
% Ricalca lo stile delle dispense di biologia scritte a mano dallo studente.
% =============================================================================
\documentclass[11pt, oneside, openany]{book}

% --- Lingua e codifica ---
\usepackage[italian]{babel}
\usepackage[T1]{fontenc}
\usepackage[utf8]{inputenc}
\usepackage{lmodern}

% --- Margini e layout (ottimizzato per lettura a schermo, non per la stampa) ---
% Margini minimi, pagina singola (niente rilegatura né margini specchiati).
\usepackage[
  margin=12mm, headsep=3mm, footskip=6mm,
  includeheadfoot
]{geometry}
\setlength{\headheight}{15pt}
\raggedbottom

% --- Spaziatura e tipografia ---
\usepackage{setspace}
\setstretch{1.15}
\setlength{\parindent}{0pt}
\setlength{\parskip}{4pt}
\usepackage[protrusion=true, expansion=true, babel=true]{microtype}

% --- Titoli capitolo (deve stare prima di hyperref) ---
\usepackage{titlesec}
\titleformat{\chapter}[block]
  {\normalfont\Huge\bfseries}
  {\thechapter\quad}{0pt}{}
\titlespacing*{\chapter}{0pt}{-60pt}{6pt}
% Intestazioni di sezione: nere, standard (niente colori).

% --- Link e segnalibri PDF ---
\usepackage[hidelinks]{hyperref}
\usepackage{bookmark}

% --- Intestazioni e piè di pagina ---
\usepackage{fancyhdr}
\pagestyle{fancy}
\fancyhf{}
\fancyhead[L]{\leftmark}
\fancyhead[R]{\rightmark}
\fancyfoot[C]{\thepage}
\renewcommand{\headrulewidth}{0.4pt}
\renewcommand{\footrulewidth}{0pt}
\fancypagestyle{plain}{%
  \fancyhf{}%
  \fancyfoot[C]{\thepage}%
  \renewcommand{\headrulewidth}{0pt}%
}
\renewcommand{\chaptermark}[1]{\markboth{Cap.~\thechapter:\ #1}{}}
\renewcommand{\sectionmark}[1]{\markright{\thesection.\ #1}}

% --- Matematica e grafica ---
\usepackage{amsmath}
\usepackage{amssymb}
\usepackage{graphicx}

% --- Tabelle ---
\usepackage{booktabs}
\usepackage{array}
\usepackage{tabularx}
\usepackage{multirow}

% --- Liste annidate (strumento principale dello stile schematico) ---
\usepackage{enumitem}
\setlist[itemize]{noitemsep, topsep=2pt}
\setlist[enumerate]{noitemsep, topsep=2pt}
% Nidificazione di default: enumerate 1. -> (a) -> i. ; itemize • -> – -> *

% --- Due colonne per liste lunghe parallele ---
\usepackage{multicol}

% =============================================================================
% COMANDI DI STILE SCHEMATICO
% =============================================================================
% Sotto-titolo in grassetto, non numerato, su riga propria (es. "Struttura", "DNA").
\newcommand{\rubrica}[1]{\par\addvspace{4pt}\noindent\textbf{#1}\par\nobreak}

% Alberi di classificazione: minimali, solo testo + rami (nessun box).
% Uso: \begin{center}\begin{forest} classalbero [Radice [A][B]] \end{forest}\end{center}
\usepackage{forest}
\forestset{
  classalbero/.style={
    for tree={
      grow'=east, parent anchor=east, child anchor=west,
      font=\footnotesize, anchor=west, align=center, edge={-},
      l sep=8mm, s sep=2mm, inner sep=1pt,
    }
  }
}

% --- Metadati ---
\title{Dispense di \underline{Doping - prevenzione e controllo}}
\author{}
\date{A.A. 2025/2026}

% =============================================================================
\begin{document}

\frontmatter                 % numerazione romana (i, ii, ...) per titolo e indice

\maketitle
\thispagestyle{empty}
\clearpage

\tableofcontents
\clearpage

\mainmatter                  % la numerazione delle pagine riparte da 1

% =============================================================================
% CAPITOLI — aggiungere \input{capitoli/nome_file} man mano
% =============================================================================

% mazzi: 01
\chapter{Le fonti del diritto internazionale in materia di doping}

% ---------------------------------------------------------------------------
\section{Il quadro normativo internazionale}

Il \textbf{legislatore internazionale} in materia di antidoping si articola su due linee, che
producono fonti distinte.

\begin{center}
\begin{forest} classalbero
[Legislatore internazionale
  [Consiglio d'Europa \\ UNESCO
    [Convenzione Strasburgo 1989]
    [Convenzione Parigi 2005]
  ]
  [WADA
    [Codice mondiale e standard]
  ]
]
\end{forest}
\end{center}

In \textit{ambito internazionale} la materia dell'antidoping è disciplinata da:
\begin{enumerate}
  \item \textbf{Convenzione contro il doping di Strasburgo}: adottata il 16 novembre 1989 e
    ratificata in Italia con la legge 29 novembre 1995, n.\,522;
  \item \textbf{Convenzione internazionale contro il doping UNESCO}: adottata a Parigi dalla
    XXXIII Conferenza generale UNESCO il 19 ottobre 2005 e ratificata in Italia con la legge 26
    novembre 2007, n.\,230;
  \item \textbf{Codice WADA} (\textit{World Anti-Doping Agency}): emanato nel 2004, aggiornato
    nel 2009 e nel 2015;
  \item \textbf{Standard internazionali WADA}.
\end{enumerate}

% ---------------------------------------------------------------------------
\section{La Convenzione di Strasburgo}

\textbf{Convenzione contro il doping di Strasburgo}: adottata il 16 novembre 1989 e ratificata
in Italia con la legge 29 novembre 1995, n.\,522. Persegue due obiettivi:
\begin{itemize}
  \item \textbf{coordinamento interno}, tra istituzioni pubbliche e sportive;
  \item \textbf{cooperazione internazionale}.
\end{itemize}

% ---------------------------------------------------------------------------
\section{La Convenzione UNESCO}

\textbf{Convenzione internazionale contro il doping UNESCO}: adottata a Parigi dalla XXXIII
Conferenza generale UNESCO il 19 ottobre 2005 e ratificata in Italia con la legge 26 novembre
2007, n.\,230; sottoscritta da 191 paesi partecipanti alla Conferenza Generale di Parigi.

\rubrica{Contenuti principali}
\begin{itemize}
  \item misure più incisive rispetto alla Convenzione di Strasburgo;
  \item riconoscimento della \textbf{leadership della WADA};
  \item lista WADA e procedure di \textbf{Esenzioni a Fini Terapeutici} allegate;
  \item \textbf{Codice WADA} come appendice alla Convenzione stessa;
  \item adozione mondiale del Codice WADA.
\end{itemize}

\rubrica{Ripartizione delle responsabilità}
\begin{itemize}
  \item la \textbf{comunità sportiva} è impegnata nei regolamenti, nei controlli antidoping,
    nonché nelle sanzioni e nelle misure disciplinari;
  \item i \textbf{governi} sono i maggiori responsabili delle misure legislative (in materia di
    tutela sanitaria), delle misure finanziarie e dei laboratori (compreso il campo della
    ricerca); coordinano inoltre le politiche e l'operato dei servizi e degli organismi pubblici
    impegnati nella lotta antidoping, e incoraggiano le organizzazioni sportive ad elaborare e
    applicare le misure appropriate di loro competenza;
  \item comunità sportiva e governi si ripartiscono la responsabilità delle \textbf{misure
    educative} e dei \textbf{programmi di informazione}.
\end{itemize}

% ---------------------------------------------------------------------------
\section{Il Codice WADA}

\textbf{Codice WADA}: insieme delle norme applicate a livello internazionale che regolamenta la
materia del doping e dell'antidoping.
\begin{itemize}
  \item armonizza le politiche, le norme e i regolamenti antidoping dei singoli paesi e
    chiarisce le responsabilità dei firmatari;
  \item obbliga all'accettazione della \textbf{Lista delle sostanze e metodi proibiti}.
\end{itemize}

\rubrica{Punti salienti}
\begin{itemize}
  \item definizione di doping \textbf{più restrittiva e più precisa};
  \item lista di sostanze;
  \item puntualizzazione per i controlli:
    \begin{itemize}
      \item ordinari, senza preavviso, in e fuori competizione;
      \item elenco mirato di atleti da sottomettere a controlli mirati (\textbf{TDP});
    \end{itemize}
  \item utilizzo di \textbf{standard internazionali} uguali:
    \begin{itemize}
      \item per tutti i laboratori;
      \item per le sostanze soggette a restrizione;
      \item per le modalità di controllo.
    \end{itemize}
\end{itemize}

\rubrica{Soggetti tenuti al rispetto del Codice}
Le regole e i principi dell'antidoping devono essere seguiti da:
\begin{itemize}
  \item \textbf{Comitati Olimpici Nazionali};
  \item \textbf{Federazioni Sportive Internazionali};
  \item \textbf{Organismi di controllo nazionali} (\textbf{NADO}).
\end{itemize}

% ---------------------------------------------------------------------------
\section{Gli standard internazionali}

Il Codice WADA rimanda a sei standard internazionali, ciascuno dedicato a un ambito specifico:
\begin{itemize}
  \item \textbf{Lista delle sostanze e metodi proibiti} (\textit{The Prohibited List});
  \item \textbf{Controlli ed investigazioni} (\textit{International Standard for Testing and
    Investigations});
  \item \textbf{Esenzioni ai fini terapeutici} (\textit{International Standard for TUE});
  \item \textbf{Laboratori} (\textit{International Standard for Laboratory});
  \item \textbf{Privacy} (\textit{International Standard for Protection of Privacy and Personal
    Information});
  \item \textbf{Compliance} (\textit{International Standard For Code Compliance By
    Signatories}).
\end{itemize}

\chapter{Le fonti del diritto nazionale in materia di doping}

% ---------------------------------------------------------------------------
\section{Il quadro normativo nazionale}

Anche il \textbf{legislatore nazionale} si articola su due linee, che producono fonti distinte.

\begin{center}
\begin{forest} classalbero
[Legislatore nazionale
  [Stato
    [L.\ 376/2000]
  ]
  [NADO
    [NSA]
  ]
]
\end{forest}
\end{center}

% ---------------------------------------------------------------------------
\section{Le Norme Sportive Antidoping (NSA)}

\textbf{Norme Sportive Antidoping} (\textbf{NSA}): documento tecnico-attuativo del Codice
Mondiale Antidoping WADA e dei relativi Standard internazionali, adottato da NADO Italia. Si
articolano in tre atti:
\begin{itemize}
  \item \textbf{Codice Sportivo Antidoping} (\textbf{CSA}): attuativo del Codice WADA;
  \item \textbf{Disciplinare dei Controlli e delle Investigazioni} (\textbf{DC-I}): attuativo
    dell'\textit{International Standard for Testing and Investigations} WADA;
  \item \textbf{Disciplinare per le Esenzioni ai Fini Terapeutici} (\textbf{D-EFT}): in
    conformità con quanto previsto dall'\textit{International Standard for Therapeutic Use
    Exemptions}.
\end{itemize}

% ---------------------------------------------------------------------------
\section{L'ordinamento statale: la legge 376/2000}

\rubrica{Sanzioni penali}
\begin{itemize}
  \item \textbf{reclusione} (3 mesi -- 3 anni) e \textbf{multa} (5--100 milioni di lire) per il
    delitto di:
    \begin{itemize}
      \item procacciamento, somministrazione, assunzione o favoreggiamento dell'utilizzo di
        farmaci o sostanze dopanti;
      \item adozione o sottoposizione a pratiche mediche dopanti;
    \end{itemize}
  \item \textbf{reclusione} (2--6 anni) e \textbf{multa} (10--150 milioni) per il delitto di
    commercio illegale di farmaci e sostanze dopanti;
  \item \textbf{circostanze aggravanti} (danni, minore, dipendenti di organismi sportivi) e
    \textbf{pene accessorie} (interdizione dall'esercizio della professione).
\end{itemize}

% ---------------------------------------------------------------------------
\section{La definizione di doping tra i due ordinamenti}

\subsection{Ordinamento statale --- legge 376/2000}

\begin{quote}
\textit{``Costituiscono doping la somministrazione o l'assunzione di farmaci o di sostanze
biologicamente o farmacologicamente attive e l'adozione o la sottoposizione a pratiche mediche
non giustificate da condizioni patologiche ed idonee a modificare le condizioni psicofisiche o
biologiche dell'organismo al fine di alterare le prestazioni agonistiche degli atleti e comunque
idonee a modificare i risultati dei controlli sull'uso dei farmaci''}
\end{quote}

\subsection{Ordinamento sportivo --- il Codice WADA}

\begin{quote}
\textit{``Per doping si intende la violazione di una o più norme antidoping definite
dall'articolo 2.1 all'articolo 2.10 del Codice''}\\
\hfill --- Art.\ 1 del Codice WADA
\end{quote}

Le Norme Sportive Antidoping elencano dieci fattispecie di violazione (artt.\ 2.1--2.10):
\begin{enumerate}
  \item \textbf{presenza} di una sostanza vietata o dei suoi metaboliti o marker nel campione
    biologico dell'atleta (Art.\ 2.1);
  \item \textbf{uso o tentato uso} di una sostanza vietata o di un metodo proibito da parte di
    un atleta (Art.\ 2.2);
  \item \textbf{elusione o rifiuto} di sottoporsi al prelievo dei campioni biologici (Art.\ 2.3);
  \item \textbf{mancata reperibilità} (Art.\ 2.4);
  \item \textbf{manomissione o tentata manomissione} in relazione a qualsiasi fase dei controlli
    antidoping (Art.\ 2.5);
  \item \textbf{possesso} di sostanze vietate e metodi proibiti (Art.\ 2.6);
  \item \textbf{traffico} illegale o tentato traffico di sostanze vietate o metodi proibiti
    (Art.\ 2.7);
  \item \textbf{somministrazione o tentata somministrazione} ad un atleta durante le
    competizioni, di un qualsiasi metodo proibito o sostanza vietata, oppure somministrazione o
    tentata somministrazione ad un atleta, fuori competizione, di un metodo o di una sostanza che
    siano proibiti fuori competizione (Art.\ 2.8);
  \item \textbf{assistenza, incoraggiamento e aiuto, istigazione, complicità} in riferimento a
    una qualsiasi violazione o tentata violazione delle NSA (Art.\ 2.9);
  \item \textbf{divieto di associazione} (Art.\ 2.10).
\end{enumerate}

% ---------------------------------------------------------------------------
\section{L'Organizzazione Nazionale Antidoping (NADO)}

\textbf{NADO}: l'ente designato da ciascun Paese quale organismo investito dell'esclusiva
autorità e responsabilità in ordine a:
\begin{itemize}
  \item adozione e attuazione delle \textbf{norme antidoping};
  \item pianificazione e conduzione dei \textbf{controlli};
  \item \textbf{gestione dei risultati} delle analisi;
  \item svolgimento dei \textbf{processi}, in ambito nazionale;
  \item irrogazione delle \textbf{sanzioni sportive}.
\end{itemize}

Secondo la normativa WADA, in Italia il \textbf{C.O.N.I.} è la \textbf{N.A.D.O.} (Organizzazione
Nazionale Antidoping).

\rubrica{Requisiti richiesti dalla WADA alla NADO}
\begin{itemize}
  \item accettazione integrale del \textbf{Programma Mondiale Antidoping WADA}:
    \begin{itemize}
      \item Codice Mondiale Antidoping;
      \item Lista Sostanze e Metodi proibiti;
      \item Standard Internazionali (controlli, laboratori, esenzioni ai fini terapeutici);
      \item competenza del \textbf{TAS di Losanna};
    \end{itemize}
  \item certificazione di un sistema di qualità \textit{``UNI ISO 9001:2000''}.
\end{itemize}

% ---------------------------------------------------------------------------
\section{Strutture e funzioni di NADO Italia}

NADO Italia è organizzata in quattro strutture, ciascuna responsabile di una funzione:
\begin{itemize}
  \item \textbf{CCA}: pianificazione dei controlli;
  \item \textbf{PNA} (\textit{Procura Nazionale Antidoping}): gestione del risultato e delle
    indagini;
  \item \textbf{CEFT}: rilascio delle esenzioni ai fini terapeutici;
  \item \textbf{TNA}: organo giudicante.
\end{itemize}

% mazzi: 03,04
\chapter{Norme sportive antidoping (NSA)}

% ---------------------------------------------------------------------------
\section{Il Codice Sportivo Antidoping}

\textbf{Codice Sportivo Antidoping} (\textbf{CSA}): attuativo del Codice WADA, disciplina la
definizione di doping e le fattispecie che ne costituiscono violazione.

% ---------------------------------------------------------------------------
\section{Articolo 1 -- Definizione di doping}

Per doping si intende la violazione di una o più norme contenute nei successivi articoli 2 e 3
del presente Codice.

% ---------------------------------------------------------------------------
\section{Articolo 2 -- Violazioni del Codice WADA}

Per doping si intende la violazione di una o più norme contenute negli articoli dal 2.1 al 2.10
del presente Codice. Lo scopo dell'articolo 2 è quello di specificare le circostanze e la
condotta che costituiscono violazione delle norme antidoping: i procedimenti nei casi di doping
saranno basati sul presupposto che una o più di queste specifiche norme siano state violate. Gli
\textbf{Atleti} o altre Persone saranno tenuti a conoscere cosa costituisca violazione delle
norme antidoping e le sostanze vietate e i metodi proibiti inclusi nella Lista.

\subsection{2.1 -- La presenza di una sostanza vietata}

\textbf{Art.\ 2.1}: la presenza di una sostanza vietata o dei suoi metaboliti o marker nel
campione biologico dell'Atleta.
\begin{itemize}
  \item \textbf{2.1.1}: ciascun Atleta deve accertarsi personalmente di non assumere alcuna
    sostanza vietata, poiché sarà ritenuto responsabile per il solo rinvenimento nei propri
    campioni biologici di qualsiasi sostanza vietata, metabolita o marker; ai fini
    dell'accertamento della violazione delle NSA non è necessario dimostrare il dolo, la colpa,
    la negligenza o l'uso consapevole da parte dell'Atleta;
  \item \textbf{2.1.2}: costituisce prova sufficiente di violazione uno dei seguenti casi:
    \begin{itemize}
      \item la presenza nel campione biologico A di una sostanza vietata o dei suoi metaboliti o
        marker, nel caso in cui l'Atleta rinunci all'analisi del campione biologico B e
        quest'ultimo non venga analizzato;
      \item la presenza nel campione biologico B di una sostanza vietata o dei suoi metaboliti o
        marker che confermi l'esito delle analisi effettuate sul campione biologico A: il
        campione biologico B dell'Atleta viene suddiviso in due flaconi e le analisi del secondo
        flacone confermano la presenza rinvenuta nel primo;
    \end{itemize}
  \item \textbf{2.1.3}: la mera presenza di un qualsiasi quantitativo di una sostanza vietata,
    dei suoi metaboliti o marker nel campione biologico dell'Atleta costituisce di per sé una
    violazione delle NSA, fatta eccezione per le sostanze per le quali la Lista delle sostanze e
    dei metodi proibiti indica specificamente un valore soglia;
  \item \textbf{2.1.4}: in deroga alla norma generale prevista dall'articolo 2.1, la Lista delle
    sostanze e dei metodi proibiti, ovvero gli Standard Internazionali, possono fissare alcuni
    criteri specifici per la valutazione delle sostanze vietate che possono essere prodotte per
    via endogena.
\end{itemize}

\subsection{2.2 -- Uso o tentato uso}

\textbf{Art.\ 2.2}: uso o tentato uso di una sostanza vietata o di un metodo proibito da parte
di un Atleta.
\begin{itemize}
  \item \textbf{2.2.1}: spetta a ogni Atleta accertarsi personalmente di non assumere alcuna
    sostanza vietata o di non utilizzare alcun metodo proibito; ai fini dell'accertamento della
    violazione delle NSA non sarà necessario dimostrare che vi sia dolo, colpa, negligenza o uso
    consapevole da parte dell'Atleta;
  \item \textbf{2.2.2}: il successo o il fallimento dell'uso o del tentato uso di una sostanza
    vietata o di un metodo proibito non costituiscono un elemento essenziale: è sufficiente,
    infatti, che la sostanza vietata o il metodo proibito siano stati usati o si sia tentato di
    usarli per integrare una violazione delle NSA.
\end{itemize}

\subsection{2.3 -- Elusione, rifiuto od omissione del prelievo}

\textbf{Art.\ 2.3}: eludere, rifiutarsi od omettere di sottoporsi al prelievo dei campioni
biologici; ovvero eludere il prelievo dei campioni biologici, oppure, senza giustificato motivo,
rifiutarsi di sottoporsi al prelievo previa notifica, in conformità alla normativa antidoping
applicabile.

\subsection{2.4 -- Mancata reperibilità}

\textbf{Art.\ 2.4}: mancata reperibilità (\textit{Whereabouts Failures}) --- ogni violazione
delle condizioni previste per gli Atleti che devono sottoporsi ai Controlli fuori competizione,
incluse la mancata presentazione di informazioni utili sulla reperibilità e la mancata esecuzione
di test in base a quanto previsto dal D-CI. Costituisce violazione ogni combinazione di 3 (tre)
mancati controlli e/o omesse comunicazioni, entro un periodo di 12 (dodici) mesi, da parte di un
Atleta inserito in un RTP.

\subsection{2.5 -- Manomissione}

\textbf{Art.\ 2.5}: manomissione o tentata manomissione in relazione a qualsiasi fase dei
Controlli antidoping --- condotta volta a minare la procedura di controllo antidoping ma che non
rientra nella definizione di pratiche vietate. La manomissione comprende, a titolo puramente
esemplificativo:
\begin{itemize}
  \item intralciare o tentare di intralciare intenzionalmente l'operato di un addetto al
    controllo antidoping;
  \item fornire informazioni fraudolente ad una Organizzazione Antidoping;
  \item intimidire o tentare di intimidire un potenziale testimone.
\end{itemize}

\subsection{2.6 -- Possesso}

\textbf{Art.\ 2.6}: possesso di sostanze vietate e ricorso a metodi proibiti.
\begin{itemize}
  \item \textbf{2.6.1}: possesso da parte di un Atleta, durante le competizioni, di qualsiasi
    sostanza vietata o il ricorso a qualsiasi metodo proibito, oppure possesso da parte di un
    Atleta, fuori competizione, di un metodo o di una sostanza espressamente vietati fuori
    competizione, a meno che l'Atleta possa dimostrare che il possesso sia dovuto ad un uso
    terapeutico consentito (artt.\ 14 e 15) o ad altro giustificato motivo;
  \item \textbf{2.6.2}: possesso da parte del Personale di supporto dell'Atleta, durante le
    competizioni, di qualsiasi sostanza vietata o di qualsiasi metodo proibito, oppure possesso
    da parte del Personale di supporto dell'Atleta, fuori competizione, di una sostanza o di un
    metodo espressamente vietati fuori competizione, in relazione a un Atleta, una competizione o
    un allenamento, a meno che il Personale possa dimostrare che il possesso sia dovuto ad un uso
    terapeutico consentito (artt.\ 14 e 15) o ad altro giustificato motivo.
\end{itemize}

\subsection{2.7 -- Traffico}

\textbf{Art.\ 2.7}: traffico illegale o tentato traffico illegale di sostanze vietate o metodi
proibiti.

\subsection{2.8 -- Somministrazione}

\textbf{Art.\ 2.8}: somministrazione o tentata somministrazione ad un Atleta durante le
competizioni, di una qualsiasi sostanza vietata o metodo proibito, oppure somministrazione o
tentata somministrazione ad un Atleta, fuori competizione, di una sostanza o di un metodo che
siano proibiti fuori competizione.

\subsection{2.9 -- Complicità}

\textbf{Art.\ 2.9}: fornire assistenza, incoraggiamento e aiuto, istigare, dissimulare o
assicurare ogni altro tipo di complicità intenzionale in riferimento a una qualsiasi violazione
o tentata violazione delle NSA o violazione dell'articolo 4.12.1 da parte di altra persona.

\subsection{2.10 -- Divieto di associazione}

\textbf{Art.\ 2.10}: l'Atleta o altra Persona soggetti all'autorità di una ADO, che si sia
avvalsa o abbia favorito la consulenza di Personale di supporto dell'Atleta:
\begin{itemize}
  \item \textbf{2.10.1}: soggetto all'autorità di una Organizzazione Antidoping che stia
    scontando un periodo di squalifica; ovvero
  \item \textbf{2.10.2}: non soggetto all'autorità di una Organizzazione Antidoping, che sia
    stato condannato o ritenuto colpevole solo nell'ambito di un procedimento penale,
    disciplinare o professionale per aver assunto una condotta che costituisca violazione del
    regolamento antidoping ove fossero stati applicati il Codice WADA e le NSA;
  \item \textbf{2.10.3}: funga da copertura o intermediario per un soggetto di cui agli artt.\
    2.10.1 o 2.10.2.
\end{itemize}

% ---------------------------------------------------------------------------
\section{Articolo 3 -- Altre violazioni delle Norme Sportive Antidoping}

Le seguenti voci costituiscono altre violazioni delle NSA:
\begin{itemize}
  \item \textbf{3.1}: qualsiasi violazione riferita alle fasi del controllo antidoping disposto
    dalla SVD di cui alla legge 376/2000;
  \item \textbf{3.2}: la mancata collaborazione da parte di qualunque soggetto per il rispetto
    delle NSA, ivi compresa l'omessa denuncia di circostanze rilevanti ai fini dell'accertamento
    di fatti di doping;
  \item \textbf{3.3}: la condotta offensiva nei confronti del DCO e/o del Personale addetto al
    controllo antidoping, la quale non sia configurabile come violazione dell'articolo 2.5 del
    CSA.
\end{itemize}

% ---------------------------------------------------------------------------
\section{Il Disciplinare per le Esenzioni ai Fini Terapeutici (D-EFT)}

\textbf{Disciplinare per le Esenzioni ai Fini Terapeutici} (\textbf{D-EFT}): attuativo
dell'\textit{International Standard for Therapeutic Use Exemption} WADA.

\textbf{Esenzione a fini terapeutici} (\textit{Therapeutic Use Exemption}, \textbf{TUE}): un
Atleta che, per motivi di salute, deve ricorrere all'assunzione e/o somministrazione di farmaci
contenenti principi attivi compresi nella Lista, deve presentare una domanda di TUE.
\begin{itemize}
  \item gli atleti di livello \textbf{Nazionale} devono presentare la domanda di TUE al
    \textbf{CEFT};
  \item gli atleti di livello \textbf{Internazionale} devono presentare la domanda di TUE al
    \textbf{TUEC} della Federazione Internazionale (FI) di appartenenza.
\end{itemize}

\subsection{Criteri per la concessione di una TUE}

\begin{itemize}
  \item \textbf{art.\ 4.1 a}: l'Atleta potrebbe subire un danno se la terapia farmacologica
    venisse sospesa;
  \item \textbf{art.\ 4.1 b}: l'uso della sostanza non deve produrre un incremento della
    performance;
  \item \textbf{art.\ 4.1 c}: non è possibile ricorrere ad una terapia farmacologica con farmaci
    non proibiti;
  \item \textbf{art.\ 4.1 d}: la necessità di utilizzare la sostanza non può essere conseguenza
    di un precedente utilizzo, senza TUE, di sostanze proibite.
\end{itemize}

\rubrica{Durata ed estinzione della TUE}
\begin{itemize}
  \item \textbf{1.2}: ciascuna TUE ha una precisa durata, così come decisa dal CEFT, al termine
    della quale cessa automaticamente di avere efficacia; se l'Atleta necessita di proseguire
    l'utilizzo della sostanza o del metodo proibiti dopo la scadenza della TUE, deve presentare,
    prima della scadenza, una nuova domanda di TUE;
  \item \textbf{1.3}: se una TUE è scaduta o è stata revocata e la sostanza proibita soggetta
    alla TUE è ancora presente nell'organismo dell'Atleta, la \textbf{Procura Nazionale
    Antidoping} (\textbf{PNA}), a seguito di un riscontro di Esito Avverso, interpella il CEFT,
    che valuta se il referto è compatibile con la scadenza o la revoca della TUE.
\end{itemize}

\subsection{Procedura per la presentazione di una domanda di TUE}

La domanda di TUE prevede la trasmissione al CEFT, a mezzo raccomandata con ricevuta di ritorno
anticipata via fax, ovvero a mezzo posta elettronica. Deve essere presentata al CEFT
\textbf{prima} di cominciare una terapia.
\begin{itemize}
  \item per le sostanze proibite in e fuori competizione, la domanda deve essere presentata
    appena formulata la diagnosi che preveda l'utilizzo di sostanze o metodi proibiti;
  \item nei casi di condizione clinica acuta e/o non procrastinabile, la domanda deve essere
    trasmessa tempestivamente \textbf{dopo} l'inizio e/o l'effettuazione (1 sola
    somministrazione) della terapia (\textbf{TUE retroattiva}).
\end{itemize}

\rubrica{Documentazione richiesta}
\begin{itemize}
  \item Modulo TUE F49 (\textit{Therapeutic Use Exemption Application});
  \item scheda per il medico curante/specialista, Modulo F51;
  \item anamnesi, storia clinica medica e documentazione comprovante la diagnosi, comprensiva dei
    risultati degli accertamenti specifici della patologia in essere e della diagnostica per
    immagini, e certificazione del medico specialista nella patologia di cui trattasi, che
    attesti sia l'assenza di eventuali controindicazioni, anche temporanee, alla pratica
    dell'attività sportiva agonistica, sia la necessità dell'utilizzo della sostanza o del metodo
    proibiti nella cura dell'Atleta, motivando le ragioni per cui non è possibile utilizzare un
    altro farmaco consentito;
  \item certificato medico di idoneità all'attività agonistica, ove previsto dalla normativa
    sportiva della disciplina di appartenenza o da norme di legge, e/o, per gli atleti
    professionisti di cui alla legge 91/1981, scheda sanitaria aggiornata con riferimento alla
    patologia per cui si richiede la TUE;
  \item certificato medico per l'attività sportiva non agonistica, ove previsto per la disciplina
    praticata.
\end{itemize}

La modulistica deve essere compilata in ogni sua parte, con redazione dattilografica o in
stampatello (``capital letter''), specificando:
\begin{itemize}
  \item Federazione Sportiva Nazionale (\textbf{FSN}) / Disciplina Sportiva Associata
    (\textbf{DSA}) / Ente di Promozione Sportiva (\textbf{EPS}) di appartenenza e la disciplina
    sportiva praticata dall'Atleta;
  \item diagnosi;
  \item principi attivi contenuti in medicinali registrati (\textit{generic name}), via di
    somministrazione (\textit{route}), dosaggio (\textit{dose}), posologia (\textit{frequency});
  \item durata di somministrazione della sostanza o dell'applicazione del metodo normalmente
    vietati per cui si richiede l'esenzione, specificando la data di inizio e la data di fine
    dell'intervento farmacologico.
\end{itemize}
La modulistica illeggibile o ritenuta incompleta non sarà esaminata e verrà restituita
all'interessato.

\subsection{Il certificato di idoneità all'attività sportiva agonistica}

Ai sensi dell'articolo 1 comma 4 della legge 376/2000, nonché delle norme per la tutela
sanitaria dell'attività sportiva agonistica contenute nei regolamenti sanitari sportivi, è cura
del medico che rilascia il certificato di idoneità informare l'Atleta in ordine agli obblighi di
conservazione di tutta la propria documentazione medica, per eventuali richieste delle Autorità
sportive.

Le esenzioni concesse dal CEFT sono comunque subordinate al rilascio, alla vigenza ovvero alla
riemissione su richiesta del CEFT stesso, del certificato di idoneità sportiva agonistica, e
comportano l'aggiornamento della scheda sanitaria per gli atleti professionisti, a norma
dell'articolo 7 della legge 23 marzo 1981, n.\,91.

\subsection{Decisione del CEFT e riconoscimento internazionale}

L'Atleta deve attendere la decisione del CEFT prima di poter fare uso della sostanza o del
metodo per cui ha richiesto la TUE.

\rubrica{Procedura e criteri di riconoscimento internazionale di una TUE}
\begin{itemize}
  \item un Atleta di livello Internazionale deve presentare la domanda di TUE alla Federazione
    Internazionale di appartenenza;
  \item laddove un Atleta di livello Internazionale abbia già una TUE concessa dal CEFT, la
    rispettiva Federazione Internazionale di appartenenza deve riconoscerne la validità.
\end{itemize}


\end{document}
